\subsection{Introduction \quad / \quad \textthai{บทนำ}}
\begin{paracol}{2}
The quest to understand the fundamental laws of the universe\index{AI Revolution@การปฏิวัติ AI (AI Revolution)} has traditionally been a journey of mathematical derivation and meticulous experimentation. From Newton's laws of motion to the intricate beauty of the Standard Model, physicists have sought to condense reality into elegant equations. However, as we venture into the era of "Big Data" physics, the complexity of our systems often outpaces our analytical capabilities. 

Enter Artificial Intelligence (AI). What was once the domain of computer science has now become a powerful lens through which we can view the physical world. This chapter explores the historical shift from traditional \textbf{Computational Physics}\index{Computational Physics@ฟิสิกส์เชิงคำนวณ (Computational Physics)} to AI-augmented science, outlining the \textbf{Paradigm Shift}\index{Paradigm Shift@การเปลี่ยนกระบวนทัศน์ (Paradigm Shift)} that is currently reshaping research at institutions like \textbf{CERN}\index{CERN@CERN (CERN)} and \textbf{LIGO}\index{LIGO@LIGO (LIGO)}.

\switchcolumn

\begin{thai}
การค้นหาความเข้าใจเกี่ยวกับกฎพื้นฐานของจักรวาลตามธรรมเนียมเดิมนั้น เป็นการเดินทางของการอนุมานทางคณิตศาสตร์และการทดลองที่พิถีพิถัน ตั้งแต่กฎการเคลื่อนที่ของนิวตันไปจนถึงความซับซ้อนที่สวยงามของ Standard Model นักฟิสิกส์มุ่งมั่นที่จะควบแน่นความเป็นจริงให้กลายเป็นสมการที่สง่างาม อย่างไรก็ตาม เมื่อเราก้าวเข้าสู่ยุคของ "ข้อมูลขนาดใหญ่ (Big Data)" ในทางฟิสิกส์ ความซับซ้อนของระบบมักจะก้าวล้ำเกินความสามารถในการวิเคราะห์แบบเดิมๆ ของเรา

การมาถึงของปัญญาประดิษฐ์ (AI) ซึ่งครั้งหนึ่งเคยเป็นขอบเขตของวิทยาการคอมพิวเตอร์ ในปัจจุบันได้กลายเป็นเลนส์อันทรงพลังที่ช่วยให้เรามองเห็นโลกทางกายภาพได้ชัดเจนขึ้น บทนี้จะสำรวจการเปลี่ยนแปลงทางประวัติศาสตร์จากฟิสิกส์เชิงคำนวณแบบดั้งเดิมไปสู่วิทยาศาสตร์ที่เสริมด้วย AI โดยเน้นย้ำถึง \textbf{การเปลี่ยนกระบวนทัศน์ (Paradigm Shift)} ที่กำลังปรับโฉมการวิจัยในสถาบันระดับโลก เช่น CERN และ LIGO
\end{thai}
\end{paracol}
